\chapter{Conclusion}
\label{ch:conclusion}

The engine does what it was built to do. Direct N-body gravity runs at
\TiledTflopsAtMaxN{} TFLOP/s of sustained single precision arithmetic at
\LadderMaxN{} particles, and the fluid solver runs \SphMaxN{} particles at
\SphMsAtMaxN{} milliseconds per step, roughly fifty times inside an interactive
budget at the fifty thousand particles it was aimed at. All \GateCount{}
physics gates pass, most with orders of magnitude of margin, against a library
compiled with fast math disabled so that no compiler flag can relax them.

The memory hierarchy ladder produced a clear answer and a clear correction.
Shared memory tiling is worth a factor of \TilingSpeedupAtMaxN{}, and the
hardware counters attribute that gain to a \SectorReduction{} times reduction
in global load sectors relieving a saturated load store pipeline, not to device
memory bandwidth, which never rose above one percent of peak in any variant.
Fusing the integrator into the force kernel was expected to be worth ten to
twenty percent and is worth \FusionGainPercentAtMaxN{} percent, for the good
reason that an $O(N^2)$ arithmetic problem on $O(N)$ data is not waiting on
memory at all. Both results are more useful than the numbers I expected to
report, because both changed how I would size the next optimisation before
writing it.

The parts of the design that survived contact with measurement did so because
they were chosen against a specific failure. Velocity Verlet was chosen because
its energy error is bounded rather than merely small, and
Figure~\ref{fig:energy} shows exactly that shape. The exact cell key comparison
in the neighbour search was added because hashing is many to one and a bucket
scan would otherwise count some neighbours twice, and the brute force
comparison confirms it works to three parts in ten million. Double precision
diagnostic accumulators were chosen because a single precision reduction would
have measured itself rather than the integrator.

\section{Further work}

A Barnes Hut tree code is the obvious next step and the one that changes the
reachable problem size rather than the constant factor. Its irregular traversal
is a genuinely different GPU problem, and the tiled kernel here would remain the
right choice for the leaf level interactions.

The fluid solver would benefit more from a position based or projection based
incompressibility treatment than from any further kernel tuning, because the
stiff equation of state, not the neighbour search, is what bounds the time step.

The tile width result suggests a small piece of unfinished business: with the
default at thirty two for fidelity to the specification, and \BestTileWidth{}
measurably faster, the right resolution is to make the width a configuration
field so the ladder can be run at the specified width while production runs use
the fastest one.
